\subsection*{一步一步建立直观}

\paragraph{1. 先明确加噪是怎样生成观测的}

右图用一个标量替代整幅图像，按照 $x_t=\sqrt{\alpha}x_0+\sqrt{1-\alpha}\,\epsilon$ 混合信号与噪声。改变保留比例，就能看出同一个观测可能主要由信号或噪声主导。这一步是已知的前向过程，还没有涉及学习如何生成图像。



\paragraph{2. 去噪依赖估计，不能凭空恢复}

若噪声估计准确，可由代数关系解出原信号；若估计偏差为 $\Delta\epsilon$，恢复误差会乘以与噪声强度有关的倍率。第二步故意用错误估计，让你比较真值与恢复值。真实模型并不知道当次噪声，需要依据训练获得的统计规律进行估计。



\paragraph{3. 为什么强噪声阶段更依赖统计结构}

当保留比例很小，恢复公式里的信号系数也很小，噪声估计误差因此被放大。图中可精确比较误差，但不代表高维去噪能唯一找回原图。扩散生成通过多步反向演化利用分布信息；这里验证的是单步代数关系，不能替代完整采样器。



\subsection*{定义、命题与证明}

\paragraph{定义：单步信号与噪声混合}

取 $0<\alpha<1$，设 $x_t=\sqrt{\alpha}x_0+\sqrt{1-\alpha}\epsilon$。给定噪声估计 $\widehat\epsilon$，按此等式解出重建值 $\widehat x_0$。

\paragraph{命题：噪声估计误差的传播}

重建误差满足 $\widehat x_0-x_0=-\sqrt{(1-\alpha)/\alpha}\,(\widehat\epsilon-\epsilon)$。

\begin{proof}

将观测的定义代入 $\widehat x_0=(x_t-\sqrt{1-\alpha}\widehat\epsilon)/\sqrt{\alpha}$，原信号项直接留下 $x_0$，剩余两项合并即可。

分母合法来自 $\alpha>0$。当 $\alpha$ 很小时，误差倍率很大；即使代数恒等式正确，也不能由它推出噪声估计足够准确或整套生成过程收敛。

\end{proof}
